\section{Uniform analytic estimates}
\label{sec:analysis}

\subsection{A coefficient bound valid on a whole disk}

\begin{lemma}[Height and divided difference]\label{lem:height}
Uniformly for \(0\le h,j\le m\),
\[
|A_{h,j,n}(\ii)|\le \exp(G_0n+o(n)).
\]
After division by \(\Pi_n\), put
\[
U_{h,n}(Y)=\sum_{j=0}^{m}
\frac{A_{h,j,n}(\ii)}{\Pi_n}Y^j\in\ZZ[\ii][Y],
\]
and define the polynomial
\[
T_{h,n}(Y)=\sum_{j=1}^{m}\frac{A_{h,j,n}(\ii)}{\Pi_n}
\sum_{a=0}^{j-1}\lambda^{j-1-a}Y^a.
\]
For \(|Y|\le2\), uniformly in \(h,Y\),
\begin{equation}\label{eq:heightbounds}
|U_{h,n}(Y)|,\ |T_{h,n}(Y)|
\le \exp(Gn+o(n)),
\end{equation}
and
\begin{equation}\label{eq:divdiff}
U_{h,n}(\lambda)-U_{h,n}(Y)
=(\lambda-Y)T_{h,n}(Y).
\end{equation}
\end{lemma}
\begin{proof}
In~\eqref{eq:polelocal}, the constant has modulus at most
\(2^{\ell_a}\). For a regular factor the constant has modulus at most
one, since it is, up to sign,
\(\ell_a!(s-1)!/(s+\ell_a)!\) for an integer \(s\ge1\).
For either type, the remaining reciprocal linear factors have
nonzero integral distances. A coefficient of fixed order at most
\(m\) is bounded by the corresponding coefficient of
\((1-w)^{-(N_n+M)}\), and is therefore polynomial in \(n\).
The factors \((-\ell+w)^h\), and the sum over \(\ell\), introduce
only more fixed powers of \(n\).
Thus
\[
|A_{h,j,n}(\ii)|
\le K\,d_n^m n^B 2^{\sum_a\ell_a}.
\]
Since \(\log d_n=n+o(n)\) and
\(\sum_a\ell_a=n\sum_j\nu_jc_j\), this has rate \(G_0\).
Division by \(\Pi_n\) subtracts \(\gamma\).
The degree \(m\) and the bounds \(|Y|\le2,\ |\lambda|<2\) are fixed,
so evaluation and the displayed finite difference introduce only
constant factors. The identity follows from
\(\lambda^j-Y^j=(\lambda-Y)\sum_{a=0}^{j-1}\lambda^{j-1-a}Y^a\).
\end{proof}

This is a uniform estimate on a disk, not just an asymptotic value
at a single point. Uniformity will permit a rational approximation
to select its own nearby algebraic point \(Y=y\).

\subsection{A global horizontal-contour bound}

The main integral estimate can be proved without leaving its tails
implicit. First observe that the solution of~\eqref{eq:rootpreview}
is unique: its left side decreases continuously from \(M\pi\) to
zero. Also
\begin{equation}\label{eq:Yhalf}
Y>\frac12.
\end{equation}
Indeed the outer type has \(c_0=1,\nu_0\ge1\), and \(M\ge2\);
at \(Y=1/2\) that type alone contributes at least \(\pi/2\), and
the remaining multiplicity makes the inequality strict.

\begin{lemma}[A uniform tail inequality]\label{lem:tail}
For \(Y>1/2\), put
\[
B_Y=\frac{(Y^2+1/4)^2}{Y^2-1/4}.
\]
If \(d=\deg Q_n=N_n+M\), then for every real \(t\),
\begin{equation}\label{eq:tailineq}
\frac{|Q_n(-n/2+nt+\ii nY)|}
 {|Q_n(-n/2+\ii nY)|}
\ge (1+t^2/B_Y)^{d/2}.
\end{equation}
The same inequality holds on the lower horizontal line.
\end{lemma}
\begin{proof}
The offsets \(u=r/n-1/2\) in each interval occur in pairs \(u,-u\),
with \(|u|\le1/2\). For a pair, the square of the modulus ratio is
\[
1+\frac{2(Y^2-u^2)t^2+t^4}{(Y^2+u^2)^2}.
\]
It is at least \((1+t^2/B_Y)^2\): the coefficient of \(t^2\) is
at least \(2/B_Y\), and that of \(t^4\) is at least \(1/B_Y^2\).
An unpaired central factor has ratio
\(\sqrt{1+t^2/Y^2}\ge\sqrt{1+t^2/B_Y}\).
Multiplication proves the claim, with repetitions included.
\end{proof}

Open the contour in~\eqref{eq:Rdef} to the lines
\(\operatorname{Im}z=\pm nY\). For each fixed \(n\) the vertical
ends vanish as their real parts tend to infinity: the rational
factor decays faster than \(1/|z|\).
On these horizontal lines, \(|\ii^{z+n}|\le e^{\Theta nY}\).
Lemma~\ref{lem:tail} therefore gives
\begin{equation}\label{eq:Rcenter}
|R_{h,n}(\ii)|
\le K_{h,Y}n^{h+1}e^{\Theta nY}
\frac{C_n}{|Q_n(-n/2+\ii nY)|}.
\end{equation}
To justify a constant independent of \(n\), the remaining integral
is bounded by
\[
\int_{\mathbb R}
\frac{(|t|+Y+1/2)^h}{(1+t^2/B_Y)^{d/2}}\dd t.
\]
For sufficiently large \(n\), replace the denominator exponent
by the fixed value \(h+2\); the resulting integral is finite.
This controls the entire line, including both infinite tails.

\subsection{Evaluating the central exponential rate}

For one interval of length \(cn\), Stirling's formula and a Riemann
sum on a smooth integrand give
\begin{align*}
\log((cn)!)&=cn\log n+cn\log c-cn+O(\log n),\\
\sum_{r=\alpha}^{\beta}
\log|r-n/2+\ii nY|
&=(cn+1)\log n\\
&\quad+n\int_{-c/2}^{c/2}\log|u+\ii Y|\dd u+O(1).
\end{align*}
The extra endpoint in the product changes only \(O(\log n)\).
Direct integration shows
\[
\int_{-c/2}^{c/2}\log|u+\ii Y|\dd u
=\frac c2\log(Y^2+c^2/4)-c
+2Y\arctan\frac{c}{2Y}.
\]
Subtracting the product logarithm from the factorial logarithm
thus gives \(-nF_c(Y)+O(\log n)\).
Equation~\eqref{eq:Rcenter} yields
\[
|R_{h,n}(\ii)|\le
\exp\left\{-n\left(\sum_j\nu_jF_{c_j}(Y)
-\Theta Y-m\right)+o(n)\right\}.
\]

\begin{proposition}[Small remainders]\label{prop:small}
After division by \(\Pi_n\), uniformly in \(0\le h\le m\),
\begin{equation}\label{eq:small}
|U_{h,n}(\lambda)|\le\exp(-Dn+o(n)).
\end{equation}
\end{proposition}

Finally,
\[
\frac{\mathrm d}{\mathrm dY}F_c(Y)
=2\arctan\frac{c}{2Y},\qquad
\frac{\mathrm d^2}{\mathrm dY^2}F_c(Y)
=-\frac{c}{Y^2+c^2/4}.
\]
Hence~\eqref{eq:rootpreview} is the unique global maximizer of
the central decay bound. This optimization needs only a
one-dimensional monotone root, not a claimed exact asymptotic
formula for an oscillating complex sum.
